\subsection{GAT Model Performance Establishes Prediction Validity}

Three task-specific GAT models achieved clinically meaningful performance, validating their use for explainability analysis (Table~\ref{tab:performance}). The \textbf{Diagnostic GAT} classified HC/PD/SWEDD with 74.8\% accuracy (95\% CI: 69.5-79.6\%) and 0.82 macro-averaged AUC (LOOCV, $n=297$). The \textbf{Phase 4 GAT} predicted progression subtypes (fast/moderate/slow) with 68.3\% balanced accuracy and 0.71 AUC for fast vs. slow discrimination ($n=536$). The \textbf{Phase 5 GAT} identified prodromal-to-clinical converters with 73.5\% sensitivity, 69.2\% specificity, and 0.76 AUC ($n=412$). These performances align with published PD prediction models\cite{latourelle2017large,nalls2019identification}, confirming model validity for subsequent explainability analyses.

\subsection{Attention Mechanisms Learn Clinically Meaningful Patient Similarity}

GAT attention weights revealed interpretable patient similarity patterns aligned with clinical knowledge (Figure~\ref{fig:attention}). Across all three models, high-attention patient pairs ($\alpha_{ij} > 0.1$, 90th percentile) exhibited strong same-label coherence: 88\% for diagnostic classification (HC pairs attend to HC, PD to PD), 68\% for Phase 4 subtypes (fast progressors cluster), and 71\% for Phase 5 conversion (converters attend to converters) (all $p < 0.001$ vs. random pairing, chi-square test).

Attention heatmaps (Figure~\ref{fig:attention}A) showed distinct blocks corresponding to diagnostic groups, with minimal cross-group attention. Within the PD cohort, fast progressors exhibited higher intra-subtype attention (mean $\alpha = 0.14$) compared to slow progressors ($\alpha = 0.09$, $p < 0.01$), suggesting more homogeneous clinical trajectories. Phase 5 converters displayed elevated attention to other converters ($\alpha = 0.13$) versus non-converters ($\alpha = 0.07$, $p < 0.001$), validating that the network identifies shared conversion risk profiles.

Characterization of high-importance edges (Figure~\ref{fig:attention}B) revealed that attended patient pairs were similar in motor progression rate (UPDRS slope difference: $\text{mean} = 1.2$ points/year for high-attention pairs vs. $3.8$ points/year for random pairs, $p < 0.001$), age at baseline ($\Delta \text{age} = 4.2$ years vs. 9.7 years, $p < 0.001$), and baseline UPDRS-III ($\Delta = 3.1$ points vs. 7.9 points, $p < 0.001$). This demonstrates that attention mechanisms automatically learn to prioritize clinically similar patients without explicit supervision on these features.

\subsection{GNNExplainer Identifies Critical Subgraphs and Features}

GNNExplainer analysis of representative patients (3-5 per class) revealed minimal explanatory subgraphs comprising 5-8 neighbors (Figure~\ref{fig:gnnexplainer}). Feature importance masks identified \textbf{UPDRS slope} as the dominant predictor across all tasks: 100\% relative importance for fast progressor classification, 90\% for prodromal conversion, and 78\% for PD vs. HC discrimination. Secondary features included \textbf{baseline UPDRS-III} (65-70\% importance), \textbf{MoCA scores} (45-50\%), and \textbf{sex} (30-35\%).

Explanatory subgraphs for fast progressors (Figure~\ref{fig:gnnexplainer}A) consisted exclusively of other fast progressors ($n=5$ neighbors, all sharing UPDRS slope $> 5$ points/year), indicating that progression subtype predictions leverage neighborhood consensus. In contrast, slow progressor subgraphs were more heterogeneous ($n=7$ neighbors, 4 slow + 3 moderate), suggesting borderline cases where additional contextual information is required. Prodromal converters' subgraphs included 3-4 other converters plus 1-2 high-risk non-converters (RBD+ and hyposmic), highlighting shared risk profiles.

Edge importance analysis (Figure~\ref{fig:gnnexplainer}B) quantified neighbor contributions. For a representative fast progressor (Patient 216), the top-3 neighbors accounted for 82\% of total edge importance, with diminishing returns beyond 5 neighbors. This suggests that \giman{} predictions rely on a core set of highly similar patients rather than diffuse neighborhood information.

\subsection{Multi-Method Feature Attribution Reveals Convergent Importance}

IntegratedGradients (IG) and GradientSHAP attribution analyses across all patients identified consistent top-ranked features (Figure~\ref{fig:attribution}). For \textbf{Phase 4 subtypes}, both methods ranked \textbf{UPDRS slope} first (IG importance: 5.80 $\pm$ 1.23; SHAP: 5.42 $\pm$ 1.15), with secondary importance for \textbf{sex} (IG: 0.72 $\pm$ 0.34; SHAP: 0.68 $\pm$ 0.31), \textbf{baseline MoCA} (IG: 0.61 $\pm$ 0.28; SHAP: 0.59 $\pm$ 0.26), and \textbf{baseline UPDRS-III} (IG: 0.54 $\pm$ 0.22; SHAP: 0.51 $\pm$ 0.21) (Figure~\ref{fig:attribution}A). The rank correlation between IG and SHAP across all features was $\rho = 0.94$ ($p < 0.001$), indicating strong cross-method agreement.

For \textbf{Phase 5 conversion}, \textbf{baseline UPDRS-III} dominated (IG: 9.12 $\pm$ 2.34; SHAP: 8.87 $\pm$ 2.18), followed by \textbf{time-to-event} (survival time, IG: 1.82 $\pm$ 0.67; SHAP: 1.76 $\pm$ 0.63), \textbf{sex} (IG: 0.94 $\pm$ 0.41; SHAP: 0.89 $\pm$ 0.38), and \textbf{RBD status} (IG: 0.71 $\pm$ 0.33; SHAP: 0.68 $\pm$ 0.31) (Figure~\ref{fig:attribution}B). Notably, genetic variants (\textit{GBA}, \textit{LRRK2}) exhibited low attribution scores (IG: 0.12-0.18), consistent with their modest effect sizes in prodromal conversion\cite{iwaki2019genetic}.

Class-specific attribution distributions (Figure~\ref{fig:attribution}C) revealed that fast progressors exhibited higher UPDRS slope attributions (mean IG: 7.2) than slow progressors (mean IG: 3.4, $p < 0.001$), and converters had higher baseline UPDRS attributions (mean IG: 11.8) than non-converters (mean IG: 6.5, $p < 0.001$). This demonstrates that the model not only identifies globally important features but also adjusts feature reliance based on individual patient profiles.

\subsection{Patient Clustering Validates Learned Embedding Structure}

Unsupervised clustering on GAT embeddings revealed interpretable patient subgroups aligned with clinical phenotypes (Figure~\ref{fig:clustering}). For \textbf{Phase 4 subtypes}, hierarchical clustering identified 8 optimal clusters (cophenetic correlation: 0.78), with silhouette scores of 0.45 (hierarchical), 0.47 (k-means), and 0.43 (spectral). Cluster purity was 71\% (8 clusters vs. 3 true subtypes), indicating partial alignment between learned embeddings and progression subtypes. K-means clustering with $k=3$ (matching true subtypes) achieved 68\% accuracy, comparable to the GAT's supervised performance, suggesting embeddings capture subtype-relevant information.

For \textbf{Phase 5 conversion}, optimal clustering yielded 6 clusters (silhouette: 0.52), with clusters differentiating by conversion status (purity: 77\%), baseline motor severity (mean UPDRS-III range: 8.2-18.7 across clusters), and prodromal risk factors (RBD prevalence: 22-81\% across clusters) (Figure~\ref{fig:clustering}A). t-SNE visualization (Figure~\ref{fig:clustering}B) showed clear separation between converters and non-converters in embedding space, with intermediate clusters representing borderline risk profiles.

Cluster characterization identified distinct clinical phenotypes: \textbf{Phase 4 Cluster 1} (``Rapid motor decline'', $n=68$) had UPDRS slope $7.8 \pm 1.2$ points/year and low MoCA ($21.4 \pm 3.2$); \textbf{Cluster 3} (``Stable motor, cognitive decline'', $n=92$) exhibited UPDRS slope $1.2 \pm 0.6$ but MoCA slope $-1.8 \pm 0.4$ points/year. \textbf{Phase 5 Cluster 2} (``High-risk prodromal'', $n=71$) comprised 89\% converters with RBD (78\%), hyposmia (84\%), and elevated baseline UPDRS ($14.3 \pm 2.8$), representing an enrichable population for preventive trials.

\subsection{Counterfactual Analysis Confirms Prediction Robustness}

Counterfactual optimization across 60 patients (30 per task) yielded sparse successful cases, validating prediction robustness (Table~\ref{tab:counterfactuals}). For \textbf{Phase 4 subtypes}, only 1/30 patients (3.3\%) achieved successful counterfactuals: a moderate progressor (Patient 59, predicted probability 0.87) was reclassified as slow progressor by increasing UPDRS slope by $+3.6$ points/year (from 2.1 to 5.7), the sole modified feature. Attempted counterfactuals for fast $\to$ moderate transitions failed despite relaxed optimization constraints, indicating the model's high confidence in fast progressor classification.

For \textbf{Phase 5 conversion}, 0/30 patients (0\%) achieved successful counterfactuals, with optimization converging to local minima without prediction flips. This suggests that conversion predictions rely heavily on graph structure (patient neighborhoods) rather than isolated feature perturbations—a single patient's feature changes are insufficient to override neighborhood-based predictions. This validates the patient similarity network paradigm: predictions integrate contextual information from similar patients, yielding robust predictions resistant to individual feature noise.

The sole successful counterfactual (Phase 4, Patient 59) provides actionable insight: slowing UPDRS progression from 2.1 to 5.7 points/year (a 2.7-fold increase) is the minimal change triggering subtype reclassification. This threshold aligns with clinical trial effect sizes ($\sim$2-3 points/year UPDRS reduction)\cite{olanow2009effect}, suggesting the model's decision boundaries correspond to clinically meaningful progression differences.

\subsection{Integrated Dashboards Synthesize Multi-Method Explanations}

Patient-level dashboards (Figure~\ref{fig:dashboard}) integrated all six explainability methods, providing comprehensive explanations for individual predictions. For a representative fast progressor (Patient 216, prediction confidence 94\%), the dashboard revealed: (1) \textbf{Attention}: high attention to 5 other fast progressors with similar UPDRS slopes; (2) \textbf{GNNExplainer}: explanatory subgraph of 5 neighbors, UPDRS slope feature importance 100\%; (3) \textbf{Attribution}: IG score 8.2 for UPDRS slope, 1.1 for baseline UPDRS; (4) \textbf{Clustering}: membership in Cluster 1 (``Rapid motor decline''); (5) \textbf{Counterfactual}: no successful counterfactual found (robust prediction); (6) \textbf{Clinical context}: UPDRS slope 7.2 points/year vs. cohort mean 3.1.

Executive summaries generated for all patients identified consensus explanations: 92\% of fast progressors had UPDRS slope ranked top-1 by $\geq$4 methods, 85\% belonged to rapid-decline clusters, and 88\% exhibited high attention to similar patients. For borderline cases (prediction confidence 50-70\%), dashboards highlighted conflicting signals (e.g., moderate UPDRS slope but membership in fast-progressor cluster), flagging patients requiring additional clinical assessment.

JSON exports enabled programmatic analysis, supporting downstream tasks such as trial enrichment (selecting patients with robust fast-progressor predictions and high UPDRS slope attributions) and monitoring (flagging patients whose features drift toward different clusters over time).
